\section{Conclusion}

\subsection{Synthèse des acquis de la première étape}

Cette analyse fonctionnelle préliminaire a permis de structurer les besoins exprimés lors de la réunion de cadrage initial autour d'une vision cohérente du système de suivi des dépenses budgétaires :

\begin{itemize}
  \item le processus métier a été organisé en quatre phases lisibles couvrant l'ensemble de la chaîne, de l'expression du besoin à la transmission du dossier d'ordonnancement au comptable public ;
  \item les cinq acteurs du système (\emph{Ordonnateur}, \emph{Fournisseur}, \emph{Contrôleur financier}, \emph{Comptable matière} et \emph{Comptable public}) ont été caractérisés avec précision, en distinguant leurs responsabilités métier réglementaires des futurs droits d'accès applicatifs ;
  \item le périmètre fonctionnel a été traduit en dix-neuf cas d'utilisation répartis en quatre domaines et représentés dans un diagramme général UML respectant strictement la notation standard.
\end{itemize}

\subsection{Prochaines étapes de conception}

Ce document de travail a désormais vocation à être examiné et validé par le responsable du projet. Les réponses apportées au questionnaire de clarification (section 8) permettront d'engager la seconde étape des travaux d'analyse et de conception, qui portera sur :

\begin{enumerate}
  \item la description textuelle détaillée des scénarios nominaux et alternatifs de chaque cas d'utilisation (notamment la gestion des rejets, des régularisations et des livraisons partielles) ;
  \item l'identification des principaux objets métier du domaine, de leurs attributs et de leurs cycles d'états (diagramme de classes du domaine métier) ;
  \item la définition de la matrice fine des habilitations par profil d'utilisateur.
\end{enumerate}
